\begin{theorem}[Charge labels derived from the regular representation]
    \label{thm:peps_regular_charge_labels}
    \lean{TNLean.PEPS.leftRegularEuclidean_adjoint,
      TNLean.PEPS.character_mem_regularChargeLabels,
      TNLean.PEPS.exists_unitary_irreducible_regularChargeLabel,
      TNLean.PEPS.regularChargeMatrix_normSq_of_mem_irreducibleCharacterFinset,
      TNLean.PEPS.regularChargeDetectorMatrix_other_of_mem_irreducibleCharacterFinset}
    \leanok
    \uses{thm:peps_unitary_character_adjoint,
      thm:peps_regular_charge_detector_matrix_eigenvectors}
    Let $G$ be a finite group and let $C$ be the finite set of irreducible
    characters occurring in its regular representation on $\ell^2(G)$.
    Every finite-dimensional irreducible complex representation has character
    in $C$. For each $\chi\in C$, an actual irreducible subrepresentation
    of $\ell^2(G)$ has character $\chi$ and inverse adjoints.
    No unitary realization of the character is supplied as a hypothesis.
    Consequently,
    \[
        \|A_\chi(p;x,y)\|^2=|G|,
        \qquad D_\chi A_\psi(p;x,y)=0\quad(\psi\in C,\ \psi\ne\chi).
    \]
    These are the group-derived charge labels and coefficient identities used
    in \cite[charge detection]{Schuch2010PEPS}, lines 2464--2486.
\end{theorem}
\begin{proof}\leanok
    The regular permutation matrices are unitary. Their invariant subspaces
    inherit inverse adjoints, so every occurring irreducible character has
    the required unitary representative. Every irreducible representation
    admits a nonzero intertwining map into the regular representation; its
    image gives an occurring irreducible subrepresentation. Apply the proved
    normalization and distinct-charge projection identities to these
    representatives.
\end{proof}
